\documentclass[../../main.tex]{subfiles}
\begin{document}
\chapter{Rings and quotient maps}
\label{ch:rings}

\section{Rings with identity}
\begin{definition}[Ring]\label{def:ring}
A ring $R$ is an abelian group under addition, together with an associative
multiplication that distributes over addition on both sides. In these notes,
rings have an identity $1$, and ring homomorphisms preserve it. The zero ring
is allowed, so the definition does not require $1\ne0$.
\end{definition}

\begin{eg}
The rings $\ZZ$, $\QQ$, $\RR$, and $\CC$ are commutative.
The ring of $2\times2$ real matrices is not commutative.
\end{eg}

\begin{prop}\label{prop:ring-zero}
For every $r\in R$, we have $r0=0r=0$.
\end{prop}
\begin{proof}
Distributivity gives $r0=r(0+0)=r0+r0$. Adding the additive inverse of
$r0$ to both sides yields $r0=0$. The other equality follows in the same way.
\end{proof}

\begin{note}
State conventions about identities and the zero ring explicitly;
they vary between authors.
\end{note}

\section{A commutative diagram}
Let $n\ge2$. The quotient homomorphism $\pi:\ZZ\to\Zc{n}$ sends an integer
to its residue class. The inclusion $\iota:n\ZZ\injto\ZZ$ has image equal
to $\Ker\pi$, and $\pi\circ\iota$ is the zero map.
\begin{marginfigure}
\centering
\begin{tikzcd}[column sep=normal,row sep=normal]
  n\ZZ \arrow[r,hook,"\iota"] \arrow[dr,"0"']
    & \ZZ \arrow[d,two heads,"\pi"] \\
    & \Zc{n}
\end{tikzcd}
\caption{The quotient map sends every multiple of $n$ to zero.}
\label{fig:quotient-map}
\end{marginfigure}
\Cref{fig:quotient-map} uses \texttt{tikz-cd}, without an external image.
For more on quotient constructions, see \textcite{judson2020}.

\section{A small multiplication table}
\begin{table}[htbp]
  \centering
  \caption{Multiplication in $\Zc{3}$. Entries are residue-class representatives.}
  \label{tab:mod-three}
  \begin{tabular}{c@{\qquad}ccc}
    \toprule
    $\cdot$ & $0$ & $1$ & $2$ \\
    \midrule
    $0$ & $0$ & $0$ & $0$ \\
    $1$ & $0$ & $1$ & $2$ \\
    $2$ & $0$ & $2$ & $1$ \\
    \bottomrule
  \end{tabular}
\end{table}
In \cref{tab:mod-three}, both nonzero elements have multiplicative inverses.
Thus $\Zc{3}$ is a field. In contrast, $\Zc{4}$ is not a field because
$2\cdot2=0$ there.

\begin{problem}
List the units of $\Zc{8}$. The macro \verb|\Zm{8}| denotes its unit group,
$\Zm{8}$.
\end{problem}
\end{document}
